\chapter{Concluzii și direcții viitoare}

\section{Rezumatul lucrării}
Lucrarea a descris proiectarea și implementarea platformei HomeForge, un sistem Smart Home cu execuție locală, conceput simultan ca aplicație funcțională și ca material didactic. Spre deosebire de soluțiile comerciale dependente de servicii cloud sau de platformele open-source mature precum Home Assistant — caracterizate de un nivel ridicat de abstractizare care reduce accesibilitatea componentelor interne pentru un dezvoltator în formare — HomeForge menține o stivă tehnologică redusă și un flux de date observabil end-to-end.

Obiectivele declarate în Capitolul 1 au fost atinse după cum urmează:
\begin{itemize}
    \item \textbf{Firmware ESP32} modular, cu patru variante template (releu, termostat cu DHT11, termostat cu BMP180, variantă combinată) și un protocol de comunicație simplu peste MQTT, descris în Capitolul 5.
    \item \textbf{Backend Django} cu un API REST complet, autentificare bazată pe JWT, sistem ierarhic de roluri \texttt{owner/admin/user} și un proces MQTT dedicat care leagă stratul hardware de baza de date PostgreSQL (Capitolul 3).
    \item \textbf{Frontend Next.js} cu dashboard interactiv, mecanism de optimistic UI pe carduri, layout persistat pe server, editor vizual pentru tipuri de dispozitive și vizualizare a topologiei prin React Flow (Capitolul 4).
    \item \textbf{Deployment containerizat} printr-o singură comandă \texttt{docker compose up}, completat cu un mod de dezvoltare nativ (Capitolul 6).
\end{itemize}

\section{Contribuții proprii}
La nivel cantitativ, contribuția proprie acoperă întreaga stivă: aproximativ 2500 de linii de Python pentru backend (\texttt{api/}), peste 12000 de linii TypeScript și JavaScript pentru frontend (\texttt{frontend/}) și cele patru proiecte de firmware ESP32 (circa 200 de linii fiecare). Pe lângă volumul de cod, câteva elemente arhitecturale particulare reprezintă contribuții substanțiale ale acestei lucrări:

\begin{enumerate}
    \item \textbf{Stocare schema-less peste RDBMS.} Starea fiecărui dispozitiv este stocată într-un \texttt{JSONField} mapat pe tipul nativ \texttt{JSONB} din PostgreSQL, în timp ce relațiile structurale (utilizator, cameră, tip de dispozitiv) rămân normalizate. Adăugarea unui nou tip de senzor nu necesită migrații SQL, iar mecanismele de control al accesului se sprijină în continuare pe integritatea referențială a schemei relaționale.
    \item \textbf{Auto-binding MAC--IP.} Mecanismul prin care procesul \texttt{mqtt\_listener} asociază automat un dispozitiv recent înregistrat (cunoscut inițial doar prin adresa IP) cu prima publicare MQTT, care conține adresa MAC. Astfel, utilizatorul nu trebuie să introducă manual adresa MAC, eliminând o sursă frecventă de erori.
    \item \textbf{Mapare „chei standard $\rightarrow$ identificatori widget".} Firmware-ul publică datele sub denumiri canonice (\texttt{temperature}, \texttt{humidity}, \texttt{relay\_1}), iar procesul de ascultare le traduce în identificatorii dinamici generați de \emph{device builder}. Drept consecință, reorganizarea widget-urilor pe card nu impune reflash-area firmware-ului.
    \item \textbf{Optimistic UI cu rollback.} Implementarea cardurilor inteligente prin \texttt{onMutate}, \texttt{onError} și \texttt{onSettled} din TanStack Query, în corespondență cu răspunsul \texttt{HTTP 202 Accepted} din backend, oferă feedback vizual imediat și revine corect la starea anterioară în cazul în care dispozitivul nu confirmă comanda. Mecanismul \texttt{matchesOptimistic} elimină situațiile în care interfața ar rămâne blocată în starea de sincronizare.
    \item \textbf{Device builder vizual.} Editorul bazat pe un graf de noduri și un panou de widget-uri, dublat de o validare server-side care impune ca fiecare \texttt{variable\_mapping} să corespundă unui ID definit în \texttt{structure[]}, permite contribuția unor tipuri noi de dispozitive fără modificări la nivelul codului platformei.
    \item \textbf{Layout cu cascadă de fallback.} Mecanismul ierarhic \emph{layout personal $\rightarrow$ layout partajat de administrator $\rightarrow$ implicit calculat client-side} permite administratorilor să stabilească o configurație inițială pentru utilizatorii noi, fără a restricționa personalizarea ulterioară a fiecăruia.
\end{enumerate}

\section{Limitări cunoscute}
Versiunea actuală a platformei prezintă o serie de limitări care trebuie semnalate explicit.

\paragraph{Securitate.} După cum s-a detaliat în Capitolul 6, brokerul Mosquitto este configurat cu \texttt{allow\_anonymous true}, comunicarea MQTT nu folosește TLS, iar variabilele Django sensibile (\texttt{SECRET\_KEY}, \texttt{DEBUG}) au valori statice. Aceste setări sunt acceptabile pentru un mediu de laborator, dar nu și pentru o instanță expusă în internet.

\paragraph{Lipsa OTA și a captive portal-ului.} Reconfigurarea sau actualizarea unui ESP32 presupune conectarea fizică prin USB. O versiune ulterioară ar putea integra ArduinoOTA sau ESPHome pentru actualizări la distanță.

\paragraph{Scalabilitate.} Cache-urile sunt locale (\texttt{LocMemCache} din Django), brokerul Mosquitto rulează în același container cu Django, iar baza de date este unică. Limita practică a implementării curente este de ordinul a câteva sute de dispozitive simultan, ceea ce acoperă scenariile rezidențiale, dar nu și pe cele de scară superioară.

\paragraph{Absența automatizărilor.} HomeForge expune dispozitivele și permite controlul manual, dar nu include încă un motor de reguli de tip \emph{trigger $\rightarrow$ condition $\rightarrow$ action}, cum oferă Home Assistant. Adăugarea acestuia constituie o direcție evidentă de continuare.

\paragraph{Testare automată.} Repository-ul conține fișierul \texttt{tests.py}, dar acoperirea testelor unitare și de integrare rămâne modestă. O eventuală evoluție către un proiect mentenabil pe termen lung impune extinderea suitei pentru a acoperi serializerele principale, fluxurile RBAC și mecanismele din \texttt{mqtt\_listener}.

\section{Direcții viitoare}
Pe baza limitărilor identificate, se pot evidenția cinci direcții de dezvoltare ulterioară.

\begin{enumerate}
    \item \textbf{TLS și autentificare MQTT.} Configurarea Mosquitto cu certificate generate automat, integrate într-un asistent de inițializare. Firmware-ul ar putea utiliza \texttt{WiFiClientSecure} cu un \texttt{rootCA} predefinit.
    \item \textbf{Motor de automatizări.} Un model Django \texttt{Automation} cu câmpurile \texttt{trigger\_topic}, \texttt{condition\_expr}, \texttt{action\_topic} și \texttt{action\_payload}, evaluat de un proces care ascultă același flux MQTT ca \texttt{mqtt\_listener}. Pe frontend, un editor vizual analog cu device builder-ul.
    \item \textbf{Suport pentru protocoale suplimentare.} Adăugarea unui adaptor pentru \texttt{Matter} sau \texttt{Zigbee2MQTT} ar permite integrarea dispozitivelor comerciale, fără abandonarea principiului \emph{local-first}. Adaptoarele ar publica pe aceleași topice \texttt{homeforge/devices/...}, astfel încât restul stivei să rămână nemodificat.
    \item \textbf{Mod public read-only.} Posibilitatea expunerii unei pagini publice (semnată cu un token) care să prezinte un subset de dispozitive în regim read-only, utilă pentru monitorizare la distanță fără a expune întreaga aplicație.
    \item \textbf{Pipeline CI/CD.} Un flux GitHub Actions care să ruleze \texttt{python manage.py test}, \texttt{npm run build} și \texttt{eslint} la fiecare \emph{pull request}, completat cu construirea imaginilor Docker la fiecare tag.
\end{enumerate}

\section{Învățări personale}
Implementarea acestei platforme a oferit ocazia parcurgerii practice a etapelor de cuplare între o stivă web modernă și un strat embedded. Patru observații se desprind ca fiind relevante pentru proiecte similare. Prima: definirea precisă a contractelor de date este mai importantă decât alegerea unui protocol anume — dezbaterea „REST vs. MQTT vs. WebSocket" devine secundară odată ce este clar ce informație circulă între componente. A doua: optimistic UI nu este o decorație vizuală, ci o strategie de gestiune a timpului perceput de utilizator, iar un mecanism de rollback robust este la fel de important ca redarea instantanee inițială. A treia: transparența arhitecturală generează un cost de implementare, dar simplifică considerabil etapa de depanare — observarea directă a traficului MQTT cu \texttt{mosquitto\_sub} s-a dovedit utilă în mod repetat. A patra: o bază de date relațională extinsă cu \texttt{JSONB} reprezintă, în numeroase scenarii, o alternativă preferabilă unei baze NoSQL pure, oferind flexibilitate fără pierderea garanțiilor ACID sau a ergonomiei unui ORM matur.

\section{Reflecție finală}
HomeForge nu își propune să acopere integralitatea funcționalităților unei platforme Smart Home comerciale. Rolul său este intermediar: face legătura între noțiunile teoretice predate la cursurile de rețele, sisteme distribuite și baze de date și aplicarea lor concretă într-un sistem care interacționează cu mediul fizic. Considerând obiectivul didactic, lucrarea poate fi utilă oricărui student sau dezvoltator care dorește să își configureze un ESP32 propriu și să urmărească, linie cu linie, traseul unei valori senzoriale de la dispozitiv până la dashboard.

Codul sursă rămâne disponibil în repository-ul proiectului, deschis utilizării, modificării și extinderii ulterioare în contextul educațional și academic.
